针对大语言模型训练中的质量评价、标度预测、预算配置和能力演进，我们按附件 A、B、C 的数据角色建立四问衔接的条件模型。第一问将 22 项质量信号定向并稳健标准化，在 A1 的 51,230 条文本上形成七域相对质量代理；对 231 对指标相关分别采用 BH-FDR 0.05 控制并在扩展集去重复核，保留 55 条复制性显著负相关边。用 A4 的 512 个配方建立 13 个目标域的二阶交互 Loss 代理，嵌套五折中 13 个目标的误差均低于线性基线；连续凸包决策的四个主要数值界间隙均低于 0.001，估算外推表则提示跨规模配方排序不稳。

第二问由 B1 的 1176 条真实轨迹拟合五参数 $N$--$D$ 主干，并用 B7 的 450 条半合成记录加入质量项。在限定共同支持域内，统一预测器显式纳入 A 侧质量代理与相对配比效应；跨附件质量坐标和 Loss 坐标无成对标定，因此该式是可检查的条件桥接。第三问在题面三项成本约束下分别求解固定配方、87 个已观测配方联立选择及原生质量敏感性；三模式各 36 格经过独立复核，$10^{19}$ FLOPs 的长上下文格明确判不可行。

第四问按模型类型及开放口径重建 C2 能力序列，以共同参数格对历史变化进行严格闭合分解；后训练组标准化均值变化 5.763 分，其中参数分布项 $-0.758$ 分、格内时间项 6.521 分。C4 资源前沿与 C2 两月 $q_{0.9}$ 前沿构成算力增速减半情景，后训练组 12 月候选为 53.828 分，条件情景范围为 37.755--72.236 分。C6 来源内 Loss 映射及 4860 行误差压力揭示部分跨题能力增益会变号。全文保留数据支持域、半合成与估算角色；纯技术因果份额、跨来源实证 Loss 换算及长期预测覆盖率均不宣称已识别。
